\section{Introduction}\hypertarget{p107.sec.1}{}
The exponent $\omega$ of matrix multiplication over $\C$ is the infimum of the real numbers $\tau$ such that, for every $\varepsilon>0$, two $n\times n$ matrices can be multiplied in $O_\varepsilon(n^{\tau+\varepsilon})$ scalar arithmetic operations. The dimension $n$ tends to infinity; the algorithm and its constants may depend on $\varepsilon$. We prove the following bound.
\begin{result}{Theorem}{1.1}{}
For every $\varepsilon>0$, two $n\times n$ complex matrices can be multiplied using $O_\varepsilon(n^{9/4+\varepsilon})$ arithmetic operations. In particular, $\omega\le9/4$.
\end{result}
Strassen's seven-product algorithm for two-by-two matrices showed that recursion can reduce the arithmetic exponent below three \cite{r13}. Bilinear algorithms can be represented by trilinear coefficient arrays, or tensors; linear substitutions and tensor products then express reductions and recursive composition. Approximate bilinear algorithms enlarge the available constructions, and Bini's interpolation argument converts them into exact asymptotic algorithms \cite{r3}. Sch\"onhage's asymptotic sum inequality turns efficient simultaneous computations of independent matrix products into bounds on the exponent \cite{r11}. Strassen's laser method extracts such independent products from tensor powers \cite{r14,r15}. His asymptotic spectrum describes asymptotic tensor comparisons through numerical invariants compatible with sums, products, and restrictions \cite{r16,r4}.

Coppersmith and Winograd combined tensor powers with progression-free sets to extract independent products from an auxiliary form \cite{r5}. Analyses of higher tensor powers by Stothers \cite{r12} and Davie--Stothers \cite{r6}, Vassilevska Williams \cite{r18}, and Le Gall \cite{r10} developed this construction; Le Gall formulated its analysis through convex optimization. Alman and Vassilevska Williams refined the extraction of independent components \cite{r2}. Duan, Wu, and Zhou analyzed losses caused by combining components across recursion levels \cite{r7}. Further asymmetric analyses by Vassilevska Williams, Xu, Xu, and Zhou \cite{r19} and Alman et al.\ \cite{r1} improved the resulting bounds. Dupont et al.\ report $\omega<2.371177$ by optimizing how components are combined across recursion levels \cite{r8}.

Our proof combines explicit tensor degenerations with Strassen's spectral viewpoint \cite{r16,r4}. Tensor characters are normalized numerical invariants that add on direct sums, multiply on tensor products, and are monotone under restriction. By \hyperlink{p107.result.2.2}{Lemma 2.2}, it suffices to bound every character on matrix multiplication. Each character has value $n^{3t}$ there for a positive parameter $t$. We bound $t$ through the growth of the character's values on polynomial multiplication. As in earlier tensor-power methods, we use interpolation and count words in which each block label occurs a prescribed number of times.

The three variable groups of a trilinear form are called legs. Useful tensor blocks can share variables, whereas a full direct sum requires independence on all three legs. Our principal construction separates blocks that are independent on two legs and share the third. We give the shared-leg variables tentative block labels, then sum substitutions with root-of-unity coefficients to cancel terms unless their indices satisfy a linear relation. We then assign integer weights to the variables so that the three weights in each surviving term sum to the square of the difference between its tentative label and its actual block label. Retaining the terms of weight zero therefore makes the labels agree. Each resulting independent block is tensored with an auxiliary dot product. When we apply this construction to an $N$th tensor power, the numbers of source copies and output blocks have a fixed ratio, independent of $N$. Its $N$th root tends to one. \hyperlink{p107.result.3.1}{Proposition 3.1} gives the construction, and \hyperlink{p107.result.3.2}{Corollary 3.2} uses it to bound character values by counting blocks with prescribed frequencies.

For polynomial multiplication with input coefficient vectors of lengths $a,b$, the output vector has length $a+b-1$. We measure this tensor by a normalized geometric mean of a character's values under all six permutations of the legs. The resulting profile $P(a,b)$ is positive and symmetric in the input lengths, and polynomial interpolation gives $P(a,b)\le(a+b-1)^{1/t}$.

Two applications of the separation inequality constrain the growth of $P$. First, tensoring with a two-dimensional dot product enlarges the second-input and output coefficient spaces. We choose a subspace of each so that multiplication by every first-input polynomial maps the input subspace into the output subspace. Multiplication on these subspaces and on their quotients gives the two neighboring input lengths. Bases adapted to these subspaces allow one choice of weights to remove the terms between the two blocks for every first input. The separation inequality then shows that successive increments of $P$ do not increase in either input length; \hyperlink{p107.sec.4.1}{Section 4.1} constructs the subspaces using a determinant.

For the second application, we partition the second-input and output coefficient indices into three matched intervals and retain only terms between matched intervals. The three resulting blocks are copies of polynomial multiplication with input lengths $a,h$, up to a permutation of the legs. Since $P$ combines all leg permutations, the separation inequality gives
\[
P(a,3h+a-1)\ge3P(a,h).
\]
Together with $P(1,b)=b$, these properties force $P(a,a)\ge a^{4/3}$. Comparing this growth with the interpolation bound yields $t\le3/4$, hence $3t\le9/4$. \hyperlink{p107.result.1.1}{Theorem 1.1} concerns the asymptotic arithmetic exponent; the proof does not specify a competitive finite matrix size.

\hyperlink{p107.sec.2}{Section 2} introduces the character framework and states the required existence result. \hyperlink{p107.sec.3}{Section 3} proves the separation inequality, and \hyperlink{p107.sec.4}{Section 4} derives the two polynomial inequalities. \hyperlink{p107.sec.5}{Section 5} proves the discrete growth estimate and completes the arithmetic argument. \hyperlink{p107.sec.A}{Appendix A} gives a complete proof of the character-existence result used here.
